% =====================================================================
% 第 VIII 章  Computation of singular j-invariants (C. Herz, Jan. 15, 1958)
% =====================================================================
\lecture{Computation of singular j-invariants}{Computation of Singular j-Invariants}{C. Herz}{Jan.\ 15, 1958}

Let $C$ be an elliptic curve. By a ``model'' of $C$ we shall mean here a
non-singular embedding in the projective plane. The first fact to be noted
is that the class of non-singular plane cubics and the class of models of
elliptic curves are identical. However there is a much deeper statement
available. The natural equivalence relation for algebraic curves is
birational equivalence, i.e.\ isomorphism of the associated function
fields; the natural equivalence relation for projective models is
projective equivalence, i.e.\ two plane curves are projectively
equivalent if and only if one may be transformed into the other by a
collineation. Projective equivalence is a priori narrower than birational
equivalence. The second main fact is that for elliptic curves the two
notions coincide provided that the ground field is algebraically closed
of characteristic $\neq 2$ or $3$. This is substantially the same as
saying that one can choose a standard cubic model for each elliptic curve
which, as it turns out, depends on a single parameter. This parameter is
essentially the modulus for the curve.

The general homogeneous cubic polynomial in three variables,
$H(x_{0},x_{1},x_{2})$, has nine coefficients and so the cubic curves
$H=0$ form an 8-dimensional projective space. If we subject the
coordinates to a linear transformation $\Pi$, a non-singular $3\times 3$
matrix, the coefficients undergo a transformation by the symmetrized
Kronecker cube of the transposed matrix:
$H^{\Pi}(x)=H(\Pi x)$ defines the transform of $H$. The group of these
transformations, obtained by taking all non-singular $\Pi$'s with
coefficients in the constant field, acts on the projective 8-space. The
homogeneous space obtained by factorization is the variety $V$ of all
plane cubic curves. It is easy to see that $H$ may be put in the form
\[
H(x_{0},x_{1},x_{2})\equiv ax_{0}x_{2}^{2}-4bx_{1}^{3}
+c_{2}x_{0}^{2}x_{1}+c_{3}x_{0}^{3}.
\]

The remaining transformations consist of the diagonal matrices $\Pi$.
Thus $V$ is the projective 3-space $(a,b,c_{2},c_{3})$ modulo the action
of the group induced by the diagonal matrices $\Pi$. This reduction step
simplifies the computations. What we are looking for are invariants of
cubic curves, -- the first things needed are the ray-invariants, i.e.\
the homogeneous polynomials, $g$, in the coefficients of the general
cubic which under the action of a matrix $\Pi$ transform by being
multiplied by a power of the determinant of $\Pi$. Any ray invariant, $g$,
is of course a ray invariant in the reduced case; the converse is true by
the ``unitary trick'' which is applicable because the general linear
group modulo the diagonal subgroup is compact. Now if $\Pi$ acting on
$\begin{pmatrix}x_{2}\\ x_{1}\\ x_{0}\end{pmatrix}$ is given by the
matrix $\begin{pmatrix}a&0&0\\ 0&\beta&\gamma\\ 0&\beta&\delta
\end{pmatrix}$, the action on $\begin{pmatrix}a&0\\ b&0\\ c_{2}&0\\
c_{3}&0\end{pmatrix}$ is given by
\[
\begin{pmatrix}
a^{2}\gamma &&&\\
&\beta^{3}&&0\\
0&&\beta\gamma^{2}&\\
&&&\gamma^{3}
\end{pmatrix}.
\]
Since this is diagonal, a basis for the ray invariants of degree $m$ is
formed by the monomials $g=a^{u}b^{v}c_{2}^{x}c_{3}^{y}$ where
$u+v+x+y=m$. Since we demand $g^{\Pi}=(\det\Pi)^{m}g$ we have the
equations $2u=m$, $3v+x=m$, $u+2x+3y=m$. The general solution is
$v=u-z$, $x=-u+3z$, $y=u-2z$ where $z$ is another integer and
$3z\geq u\geq 2z$. Taking $z=1$ there are two possibilities:
$u=2$ corresponding to the monomial $g_{2}=a^{2}bc_{2}$ and $u=3$
corresponding to $g_{3}=a^{3}b^{2}c_{3}$. It follows that all other ray
invariants are polynomials in $g_{2}$ and $g_{3}$.
\[
\Delta=g_{2}^{3}-27g_{3}^{2}=a^{6}b^{3}(c_{2}^{3}-27bc_{3}^{2})
\]
has the virtue that $\Delta=0$ is a necessary and sufficient condition
for the cubic curve to be singular, i.e.\ either have a linear factor or
a double point. $j=2^{6}\cdot 3^{3}\,g_{2}^{3}\Delta^{-1}$ is defined for
all non-singular cubics. If $V_{o}$ is the sub-variety of $V$ defined by
$\Delta=0$ then it is easy to see that $V-V_{o}$ is an affine line and
$j$ is a parameter. $V-V_{o}$ is the variety of moduli of elliptic
curves.

$\Delta$ has another virtue. Suppose $H(x_{0},x_{1},x_{2})=0$ is an
elliptic curve, $C$. Then
\[
du=2\left|\begin{array}{ccc}
\lambda_{o}&x_{o}&dx_{o}\\
\lambda_{1}&x_{1}&dx_{1}\\
\lambda_{2}&x_{2}&dx_{2}
\end{array}\right|
\Bigg/\sum_{j=o}^{2}\lambda_{j}\,\partial H/\partial x_{j}
\]
is a differential of the first kind on $C$, and is independent of the
choice of the point $(\lambda_{2},\lambda_{1},\lambda_{0})$ in the
projective plane. Let us examine the effect of a change of coordinates on
$du$. Suppose $x=\Pi x^{*}$. Then the new equation is
$H^{*}(x^{*})=H(\Pi x^{*})=H^{\Pi}(x^{*})$. In the new system
\[
du^{*}=2\left|\begin{array}{ccc}
\lambda_{o}^{*}&x_{o}^{*}&dx_{o}^{*}\\
\lambda_{1}^{*}&x_{1}^{*}&dx_{1}^{*}\\
\lambda_{2}^{*}&x_{2}^{*}&dx_{2}^{*}
\end{array}\right|
\Bigg/\sum_{j=o}^{2}\lambda_{j}^{*}\,\partial H^{*}/\partial x_{j}^{*}.
\]
Of course $du^{*}=\mu\,du$ where $\mu$ is a constant multiplier since
differentials of the first kind differ only by a constant factor. However
since $\lambda$ is arbitrary we can put $\lambda=\Pi\lambda^{*}$ and it
becomes evident that $du=(\det\Pi)\,du^{*}$. Hence $\mu=\det\Pi$ and
$\mu^{12}=\Delta^{*}/\Delta$, -- for this reason expressions of the form
$\Delta^{*}/\Delta$ are occasionally called multipliers.
$\Delta\,du^{12}$ is an invariant of the curve. Could we choose
$\Delta^{1/12}$ invariantly, $\Delta^{1/12}\,du$ would be a differential
of the first kind canonically attached to the curve.

The group operation on an elliptic curve has a simple geometric
interpretation. Take a model and let $(\mathcal{U})$ and $(\mathcal{V})$
be two points on it. These points determine a line which intersects the
cubic in exactly one other point which we write as
$(-\mathcal{U}-\mathcal{V})$. The curve has a unique point of inflection
$(\mathcal{O})$. It is easy to check that the operations defined give
rise to an additive abelian group with identity element
$(\mathcal{O})$. Since the group operation is defined by lines it is
clearly invariant under collineations.

There are exactly three points on $C$, other than $(\mathcal{O})$, whose
tangents pass through $(\mathcal{O})$. If $(\eta)$ is such a point,
$(2\eta)=(\mathcal{O})$, according to the description of the addition
operation, and the converse is true. Thus if $(\eta_{1})$ and
$(\eta_{2})$ are two such points, $(\eta_{3})=(-\eta_{1}-\eta_{2})$ is a
third. These three points are collinear; call the line they determine
$L_{C}$. Projection from $(\mathcal{O})$ gives a two-fold covering of
$C$ onto $L_{C}$; two points on $C$ go into the same point of $L_{C}$ if
and only if they are inverses of one another. (The projection is not
defined at $(\mathcal{O})$ so we extend by continuity, -- $(\mathcal{O})$
is projected along the inflection tangent.) Let $\eta_{0}$ be the
projection of $\mathcal{O}$ on $L_{C}$. The cross ratio
$(\eta_{0}\eta_{1}\eta_{2}\eta_{3})=k^{2}$ is invariant under projective
transformations; however the three points $(\eta_{1}),(\eta_{2}),
(\eta_{3})$ may be permuted at will so that $k^{2}$ is an invariant of
$C$ only up to the action of the symmetric group on three letters. If we
write $k'^{2}=1-k^{2}$, the full invariant is
$j=2^{8}\frac{(1-k^{2}k'^{2})^{3}}{k^{4}k'^{4}}$.

Suppose $C^{*}$ is an elliptic curve which forms an m-fold covering of
$C$. Since the covering map must be a homomorphism of the group
structures with suitable base points, inverses are mapped into inverses.
Thus if we have a model of $C^{*}$ with inflection point
$(\mathcal{O}^{*})$ mapping onto a model of $C$ with inflection point
$(\mathcal{O})$, there is a well-defined map of $L_{C^{*}}$ onto $L_{C}$
which is an m-fold covering of one projective line by another. It will
be convenient to take the models in the projective plane with
$(\mathcal{O})$ at $x_{0}=x_{1}=0$, $x_{2}=1$, inflection tangent
$x_{0}=0$, and the line $L_{C}$ as $x_{2}=0$. The covering map of
$L_{C^{*}}$ onto $L_{C}$ then has the form
\[
x_{0}=\mathbb{Q}(x_{0}^{*},x_{1}^{*})\,x_{0}^{*},\qquad
x_{1}=\mathbb{P}(x_{0}^{*},x_{1}^{*})
\]
where $\mathbb{Q}$ is a homogeneous polynomial of degree $m-1$ and
$\mathbb{P}$ is one of degree $m$. The covering map of $C^{*}$ onto $C$
may then be described at all but a finite number of points by the
additional transformation $x_{2}=R(x_{0}^{*},x_{1}^{*})x_{2}^{*}$ where
$R$ is a homogeneous rational function of degree $n-1$. We shall now
find the explicit form of $R$. The equations of $C$ and $C^{*}$ are
$x_{0}x_{2}^{2}-T(x_{0},x_{1})=0$ and
$x_{0}^{*}x_{2}^{*2}-T^{*}(x_{0}^{*},x_{1}^{*})=0$ respectively where $T$
and $T^{*}$ are homogeneous cubics with the coefficient of $x_{1}^{3}$
different from $0$. The differentials of the first kind are
\[
du=\frac{\left|\begin{array}{cc}x_{o}&dx_{o}\\ x_{1}&dx_{1}\end{array}\right|}
{x_{o}x_{2}}
\qquad\text{and}\qquad
du^{*}=\frac{\left|\begin{array}{cc}x_{o}^{*}&dx_{o}^{*}\\
x_{1}^{*}&dx_{1}^{*}\end{array}\right|}{x_{o}^{*}x_{2}^{*}}.
\]
Since $x_{0}$, $x_{1}$ are homogeneous polynomials of degree $m$ in
$x_{0}^{*}$, $x_{1}^{*}$,
\[
\left|\begin{array}{cc}x_{o}&dx_{o}\\ x_{1}&dx_{1}\end{array}\right|
=\frac{1}{m}\left|\;\overline{\frac{\partial(x_{o},x_{1})}
{\partial(x_{o}^{*},x_{1}^{*})}}\;\right|\cdot
\left|\begin{array}{cc}x_{o}^{*}&dx_{o}^{*}\\ x_{1}^{*}&dx_{1}^{*}
\end{array}\right|.
\]
Hence
$du=\frac{1}{m}\bigl|\frac{\partial(x_{0},x_{1})}{\partial(x_{0}^{*},
x_{1}^{*})}\bigr|\frac{du^{*}}{\mathbb{Q}R}$, but we must also have
$du^{*}=\mu\,du$ where $\mu$ is a constant multiplier. This tells us that
$R=\frac{\mu}{m}\mathbb{Q}^{-1}\bigl|\frac{\partial(x_{0},x_{1})}
{\partial(x_{0}^{*},x_{1}^{*})}\bigr|
=\frac{\mu}{m}\Bigl\{\frac{\partial\mathbb{P}}{\partial x_{1}^{*}}
+\mathbb{Q}^{-1}x_{0}^{*}\Bigl|\frac{\partial(\mathbb{Q},\mathbb{P})}
{\partial(x_{0}^{*},x_{1}^{*})}\Bigr|\Bigr\}$. Substituting in the
equation for the curves we find that
\[
\frac{\mu^{2}}{m^{2}}\Bigl\{
\mathbb{Q}\Bigl(\frac{\partial\mathbb{P}}{\partial x_{1}^{*}}\Bigr)^{2}
+2x_{0}^{*}\Bigl|\frac{\partial(\mathbb{Q},\mathbb{P})}
{\partial(x_{0}^{*},x_{1}^{*})}\Bigr|
+\mathbb{Q}^{-1}x_{0}^{*2}
\Bigl|\frac{\partial(\mathbb{Q},\mathbb{P})}
{\partial(x_{0}^{*},x_{1}^{*})}\Bigr|^{2}\Bigr\}\cdot
T^{*}(x_{0}^{*},x_{1}^{*})
=T\bigl(x_{0}^{*}\mathbb{Q},\ \mathbb{P}\bigr).
\]

There can be no denominator arising from the $\mathbb{Q}^{-1}$ term on
the left. Hence the only possible linear factors of $\mathbb{Q}$ are
$x_{0}^{*}$ and the factors of $T^{*}$; all other factors of
$\mathbb{Q}$ must be at least quadratic.

What we have just described is the Jacobi transformation principle. It
leads to the algebraic computation of the singular class invariants $j$.
Some explicit computations are performed in Weber: \emph{Lehrbuch der
Algebra}, Vol.\ III, sections 8--10. The results, of course, coincide
with the invariant equations for $j$ obtained by analytic means. These
equations will hold in any of the fields we have considered, i.e.\ the
analytic method gives results valid over rather general ground fields.
The difficulty in making computations from the class equations is due to
the fact that the coefficients in the q-expansion of $j$ are intractably
large:
\[
j=q^{-1}+744+196\,884\,q+21\,493\,760\,q^{2}+\dots
\]
It is in fact easier to compute using $k^{2}$. However we shall now
proceed to the analytic case and examine the class equation by
investigating the subgroups of finite index in the modular group.

Let $\Gamma$ denote the full modular group and $H$ a (not necessarily
normal) subgroup of index $\psi$. We shall write
$\Gamma=\sum_{a=1}^{\psi}HV_{a}$. Let $\Omega$ be the standard
fundamental domain for the modular group; $\Omega$ has a simplicial
decomposition with 2 faces, 3 edges, and 3 vertices, namely
$\tau=i\infty$, $i$, and $\rho$. $S$ and $T$ stand for the
transformations $S:\tau\longrightarrow\tau+1$, $T:\tau\longrightarrow
-1/\tau$. $i\infty$ is a fixed point of $S$ and all its powers, $i$ is a
fixed point of $T$ ($T^{2}=I$), $\rho$ is a fixed point of $S^{-1}T$
($(S^{-1}T)^{3}=I$). $\Gamma$ is generated by $S$ and $T$ and the only
relations are consequences of those given above.

A fundamental domain, $\mathfrak{F}$, for the subgroup $H$ is defined as
follows. At each interior point $\tau\in\Omega$ we assign $\psi$ points
$V_{1}\tau,\ldots,V_{\psi}\tau$ in $\mathfrak{F}$. These points are
inequivalent modulo $H$ for if $V_{a}\tau=UV_{b}\tau$ then $\tau$ is a
fixed point of $V_{a}^{-1}UV_{b}$. This can occur for
$\tau\in\Omega$, $\tau\neq i\infty,i,\rho$ only if $V_{a}^{-1}UV_{b}=I$.
Hence if $U\in H$, $V_{a}=UV_{b}$ contrary to the assumption that for
$a\neq b$, $V_{a}$ and $V_{b}$ represent different cosets. Thus
$\mathfrak{F}$ has $2\psi$ faces and $3\psi$ edges.

The vertices require special treatment. First we take $\tau=i\infty$.
Let us define an equivalence relation $V_{a}\sim V_{b}$ if and only if
$V_{a}S^{k}V_{b}^{-1}\in H$ for some integer $k$. This provides a
grouping of the cosets of $\Gamma$ mod $H$ into equivalence classes
$K_{1},\ldots,K_{d}$ with $K_{c}/H$ of order $n_{c}$. For each class
$K_{c}$ there is one vertex of $\mathfrak{F}$ above $i\infty$ in $\Omega$
at which $n_{c}$ sheets come together.

A similar but simpler analysis holds at $\tau=i$ and $\tau=\rho$.
$V_{a}$ and $V_{a}T$ represent two different sheets of $\mathfrak{F}$
unless $V_{a}TV_{a}^{-1}\in H$; let $\varepsilon_{i}$ be the number of
solutions of this equation. Then there are
$\frac{1}{2}(\psi-\varepsilon_{i})$ vertices of $\mathfrak{F}$ above $i$
in $\Omega$ at which 2 sheets come together and $\varepsilon_{i}$
unbranched vertices. Likewise $V_{a}$, $V_{a}(S^{-1}T)$, and
$V_{a}(S^{-1}T)^{2}$ represent three distinct sheets of $\mathfrak{F}$
unless $V_{a}S^{-1}TV_{a}^{-1}\in H$ in which case all three coincide;
let $\varepsilon_{\rho}$ be the number of solutions of the last equation.
There are $\frac{1}{3}(\psi-\varepsilon_{\rho})$ vertices of
$\mathfrak{F}$ above $\rho$ in $\Omega$ at which 3 sheets come together
and $\varepsilon_{\rho}$ unbranched vertices.

The total number of vertices of $\mathfrak{F}$ is
$d+\frac{1}{2}(\psi+\varepsilon_{i})+\frac{1}{3}(\psi+2\varepsilon_{\rho})$.
Hence, in particular, $\mathfrak{F}$ is a sphere if and only if
$2\psi-3\psi+d+\frac{1}{2}(\psi+\varepsilon_{i})
+\frac{1}{3}(\psi+2\varepsilon_{\rho})=2$, i.e.
$\psi=6d+3\varepsilon_{i}+4\varepsilon_{\rho}-12$.

The most important subgroups for our purposes are the transformation
groups $H_{m}$ and the congruence groups $\Gamma_{m}$. The transformation
groups arise as follows. Let $M$ be the matrix
$\begin{pmatrix}m&0\\ 0&1\end{pmatrix}$. $U\in H_{m}$ if $U\in\Gamma$ and
$MUM^{-1}\in\Gamma$ which is equivalent to saying that $U$ has the matrix
representation $U=\begin{pmatrix}\alpha&\beta\\ \gamma&\delta\end{pmatrix}$
with $\alpha,\beta,\gamma,\delta$ integers, $\alpha\delta-\beta\gamma=1$
and $\gamma\equiv 0\bmod m$. Writing
$\Gamma=\sum_{a=1}^{\psi(m)}H_{m}V_{a}$ the matrices $MV_{a}$ can be
chosen to run through the set
$\begin{pmatrix}\alpha&\beta\\ 0&\delta\end{pmatrix}$ where
$\alpha\geq 1$, $0\leq\beta<\delta$, $\alpha\delta=m$ and
$(\alpha,\beta,\delta)=1$. We shall write $V_{\alpha,\beta,\delta}$ for
$V_{a}$. The invariant polynomial is
$J(t,j)=\prod_{\alpha,\beta,\delta}\bigl(t-j(MV_{\alpha,\beta,\delta})
\bigr)$. $j^{M}$ defined by $j^{M}(\tau)=j(m\tau)$ is invariant under the
action of $H_{m}$. Thus $j^{M}$ is a meromorphic function on the Riemann
surface $\mathfrak{F}_{m}$ corresponding to $H_{m}$. The Galois group of
the invariant equation $J(t,j)=0$ over the field $\mathbb{Q}(j)$ is
$\Gamma/\Gamma_{m}$ where $\Gamma_{m}$ is the largest subgroup of $H_{m}$
which is normal in $\Gamma$. $\Gamma_{m}$ is represented by the matrices
$\begin{pmatrix}\alpha&\beta\\ \gamma&\delta\end{pmatrix}\equiv
\begin{pmatrix}1&0\\ 0&1\end{pmatrix}\bmod m$ with $\alpha,\beta,\gamma,
\delta$ integers and $\alpha\delta-\beta\gamma=1$. If $m$ is a prime
$\Gamma/\Gamma_{m}$ is isomorphic to $\mathrm{SL}_{2}(\mathbb{Z}_{m})$
modulo its center; this group is simple for $m\geq 5$.

The computations involving singular invariants $j$ are easy to carry out
when the Riemann surface $\mathfrak{F}_{m}$ corresponding to $H_{m}$ is a
sphere. Here $\psi=\psi(m)=m\prod_{p\mid m}\bigl(1+\frac{1}{p}\bigr)$.
We have to calculate $d$, $\varepsilon_{i}$, and $\varepsilon_{\rho}$.
The equivalence relation at $i\infty$,
$V_{\alpha,\beta,\delta}\sim V_{\alpha',\beta',\delta'}$ may be formulated
as $(\alpha_{0},\beta_{0})\bigl(\begin{smallmatrix}1&k\\ 0&1\end{smallmatrix}
\bigr)=U(\alpha_{0}',\beta_{0}')$ where $U\in\Gamma$. The matrix on the
left is $(\alpha_{0},\beta_{0}+k\alpha_{0})$ and the equivalence holds
only if $\alpha'=\alpha$, $\delta'=\delta$ and
$\beta'\equiv\beta+k\alpha\bmod\delta$. For simplicity assume $m$ is
square-free. Then $(\alpha,\delta)=1$ and $\beta'$ can run through all
the values $0,1,\ldots,\delta-1$ for any given $\beta$. Hence for each
$\delta\mid m$ we have an equivalence class $K_{\delta}$ containing
$\delta$ cosets. The number $d$ is $d(m)$, the number of divisors of
$m$.

The equation $V_{\alpha,\beta,\delta}TV_{\alpha,\beta,\delta}^{-1}\in
H_{m}$ is equivalent to $\delta=m$ and $\beta^{2}+1\equiv 0\bmod m$. The
analogous equation at the vertices above $\tau=\rho$ is equivalent to
$\delta=m$ and $\beta^{2}-\beta+1\equiv 0\bmod m$. For simplicity,
suppose from now on that $m$ is a prime $>3$. Then $\psi(m)=m+1$, $d=2$,
$\varepsilon_{i}=2$ or $0$ according to whether $-1$ is a quadratic
residue mod $m$ or not, $\varepsilon_{\rho}=2$ or $0$ according to
whether $-3$ is a quadratic residue mod $m$ or not. The condition for
$\mathfrak{F}_{m}$ to be a sphere is $m+1=3\varepsilon_{i}
+4\varepsilon_{\rho}$. This holds for $m=5,7,13$ and no other primes $>3$.

We shall now carry out an extended analysis for the case $m=5$. On
$\mathfrak{F}_{5}$ there is one vertex above $\tau=i\infty$ at which
there is only one sheet and one vertex at which 5 sheets come together.
Above $\tau=i$ there are two vertices at which 2 sheets come together and
two unbranched vertices. Above $\tau=\rho$ there are two vertices at
which three sheets come together.

We know in advance that
$\Delta\bigl(\tfrac{1}{5}\omega_{1},\omega_{2}\bigr)
/\Delta(\omega_{1},\omega_{2})$ is a meromorphic function on
$\mathfrak{F}_{m}$ since it is invariant under $H_{m}$. The only possible
zeros and poles occur at the vertices lying over $\tau=i\infty$. Let $q$
be the local uniformizing parameter $e^{2\pi i\omega_{2}/\omega_{1}}$ at
the unbranched vertex over $\tau=i\infty$. Here
\[
\frac{\Delta(\frac{1}{5}\omega_{1},\omega_{2})}{\Delta(\omega_{1},\omega_{2})}
=5^{12}\,q^{4}\prod_{5\nmid n}(1-q^{n})^{-24}.
\]

The function has a four-fold zero; since it has only one pole, the pole
is also 4-fold and this function is the fourth power of a function $f$
with only one single pole. Since $\mathfrak{F}_{5}$ is a sphere, $f$ is a
global uniformizing parameter. Now consider the function $j$. $j$ has a
six-fold pole at the branched vertex above $\tau=i\infty$. At each of the
two vertices above $\tau=\rho$ it has a three-fold zero. Hence
$j=v^{3}$ where $v$ is a meromorphic function on $\mathfrak{F}_{5}$.
Since $v$ has a double pole at the pole of $f$, $v$ is in fact a
quadratic polynomial in $f$. Let
\[
q'=e^{-2\pi i\frac{\omega_{1}}{5\omega_{2}}}
\]
be the local uniformizing parameter at the branched vertex above
$\tau=i\infty$. At this point $j=q'^{-5}+E$ where $E$ is a power series
in non-negative powers of $q'$. $f$ is defined by
\[
f=\left\{\frac{\Delta(-\omega_{2},\frac{1}{5}\omega_{1})}
{\Delta(-\omega_{2},\omega_{1})}\right\}^{1/4}
=q'^{-1}\prod_{5\nmid n}(1-q'^{n})^{6}
\]
so that $f=q'^{-1}+\ldots$, $f^{2}=q'^{-2}-12q'^{-1}+\ldots$, and
$v=q'^{-2}-2q'^{-1}+\dots$. Hence $v-(f^{2}+10f)$ has no poles; it is
thus a constant. At the other vertex $j=q^{-1}+\ldots$, $f=5^{3}q+\dots$
so $v=5+\dots$ and the constant term is $5$. We have computed
$v=f^{2}+10f+5$, $j=(f^{2}+10f+5)^{3}/f$.

We now make use of a very important property of the transformation
groups: namely $(TM)H_{m}(TM)^{-1}=H_{m}$. Thus $f^{(TM)}$ is a
meromorphic function on $\mathfrak{F}_{m}$ if $f$ is. In the case $m=5$
let us write $\tilde{f}=f^{(TM)}$, i.e.\ $\tilde{f}(\tau)
=f(-\frac{1}{5\tau})$. The q-expansion of $\tilde{f}$ is the same as the
$q'$-expansion of $f$, i.e.\ $\tilde{f}=q^{-1}+\dots$. Since
$\tilde{\tilde{f}}=f$, $\tilde{f}$ has only one pole, and we conclude
that $\tilde{f}=5^{3}f^{-1}$. For $\tau=\sqrt{-5}/5$,
$-\frac{1}{5\tau}=\tau$ so $f=\tilde{f}$ and hence
$f(\sqrt{-5}/5)=\pm 5^{3/2}$. With either choice of sign
$f^{2}+10f+5>0$; hence the sign is the same as that of
$j(\sqrt{-5}/5)$, but $j$ is positive on the imaginary axis. Accordingly
$f(\sqrt{-5}/5)=+5^{3/2}$. $(5,\sqrt{-5})$ is an ideal basis in the field
$\mathbb{Q}(\sqrt{-5})$ representing the principal class. From the
expression of $j$ in terms of $f$ one has
$j=2^{3}\cdot(25+13\sqrt{5})^{3}$ for the singular invariant of
the principal class.
% 校对注记：由 $j=(f^{2}+10f+5)^{3}/f$ 与 $f(\sqrt{-5}/5)=5^{3/2}$ 得
% $j=2^{3}(25+13\sqrt{5})^{3}=632000+282880\sqrt{5}$（数值与经典值一致；
% 原书末行排印作 $2^{3}\cdot 5\cdot(25+13\sqrt{5})^{3}$，应为衍文）。
